\begin{abstract}
This note takes stock of a representation-aware constructive programme for
four-dimensional Yang--Mills theory.  The programme has produced exact local
and finite-block results, counterexamples to scalar worst-history estimates,
and conditional representation-to-spectrum theorems, but it has not proved
physical one-scale basin entry, repeated renormalization-group control,
continuum Yang--Mills theory, or a Hamiltonian mass gap.  We argue that the
missing object is a normalized, state-weighted, topologically complete
transfer dynamics between scales.  We formulate a minimal substrate as a
region--scale network of operator systems, states, and completely positive
maps; state ten aggressive hypotheses; and report executable first probes of
each.  Two probes materially sharpen the programme: coordinated multi-star
constraints acquire zero-spill directions on the certified finite block, and
center labels together with a Pfaffian charge fail to compose under addition
of a common flux context.  A reward-ranked strategy selects a
vacuum-weighted, normalized, multichannel one-step kernel inequality as the
primary target, with coordinated polymer cancellation as the leading backup.
All continuum and spectral conclusions remain explicitly conditional.
\end{abstract}
